
On
définit ensuite un homomorphisme \(K_1\times K'\to K_1\) en associant à \((y',x)\) l'élément \(y'(1-L)+g^*\circ i^*(x)\); cet homomorphisme s'annule sur \(K\), et définit par conséquent un homomorphisme \(j^*:K'_1\longrightarrow K_1\). Enfin, on définit un homomorphisme \(K_1\times K'\to K'\) en associant à \((y',x)\) l'élément \(i_*\circ g_*(y')+x\); cet homomorphisme s'annule sur \(K\) et définit donc un homomorphisme \(f_*:K'_1\longrightarrow K'\).

\textbf{Proposition 4.7.6.}
\begin{enumerate}[label=\roman*)]
\item \(j^*\) \uline{et} \(f^*\) \uline{sont des} \(\lambda\)-\uline{homomorphismes};
\item \(j_*\) \uline{est} \(K'_1\)-\uline{linéaire et} \(f_*\) \uline{est} \(K'\)-\uline{linéaire};
\item \(j^*\circ j_*\) \uline{est l'homomorphisme} « \uline{multiplication par} \(1-L\) »;
\item \uline{pour tout} \(y'\in K_1\), \uline{et tout} \(p\ge1\), \(j_*\bigl(\lambda^p(L,y')\bigr)=\lambda^p(j_*(y'))\);
\item \(f_*\circ f^*=\Id\).
\end{enumerate}

Il suffit de faire tourner la machine.

Si on en revient à la situation géométrique rappelée plus haut, on voit que \(K^{\bullet}(X')\) se construit à partir de \(K^{\bullet}(Y)\), \(K^{\bullet}(X)\), des applications \(i^*\) et \(i_*\) et de l'élément \(N\) de \(K^{\bullet}(Y)\) par la méthode précédente, les opérations \(\lambda\) sur \(K^{\bullet}(X')\) ne coïncidant toutefois avec les opérations \(\lambda\) définies plus haut sur \(K'_1\) qu'à torsion près, tant que la relation (4.3.1) n'est prouvée que modulo torsion.

\textbf{Théorème 4.8.} \uline{Soient} \(i:Y\to X\) \uline{une immersion fermée régulière de codimension constante} \(d\), \(X'\) \uline{le schéma obtenu en faisant éclater} \(Y\), \uline{et} \(Y'=X'\times_XY\), \uline{de sorte qu'on a le diagramme commutatif} :
\[
\begin{tikzcd}[column sep=large,row sep=large]
Y' \arrow[r,"j"] \arrow[d,"g"'] & X' \arrow[d,"f"]\\
Y \arrow[r,"i"] & X.
\end{tikzcd}
\]
\uline{On suppose que} \(X\) \uline{est quasi-compact et a un faisceau ample}. \uline{Alors on a pour tout} \(k\) \uline{la suite
exacte} :
